\documentclass[a4paper,10.5pt]{article}
\usepackage[utf8]{inputenc}
\usepackage[T1]{fontenc}
\usepackage[french]{babel}
\usepackage[left=1cm,right=1cm,top=.5cm,bottom=0cm]{geometry}
\usepackage{enumitem}
\usepackage{hyperref}
\usepackage{xcolor}
\usepackage{titlesec}
\usepackage{fontawesome5}
\usepackage{lmodern}
\usepackage[normalem]{ulem}

% --- Couleurs ---
\definecolor{primary}{RGB}{35, 55, 123}
\definecolor{secondary}{RGB}{80, 80, 80}

% --- Configuration des liens ---
\hypersetup{colorlinks=true, linkcolor=primary, filecolor=primary, urlcolor=primary}

% --- Formatage des titres ---
\titleformat{\section}{\large\bfseries\color{primary}\uppercase}{}{0em}{}[\titlerule]
\titlespacing{\section}{0pt}{8pt}{5pt}

% --- Commandes personnalisées ---
\newcommand{\resumeItem}[1]{\item\small{#1}}
\newcommand{\resumeSubheading}[4]{
  \vspace{-1pt}\item
    \begin{tabular*}{0.97\textwidth}[t]{l@{\extracolsep{\fill}}r}
      \textbf{#1} & \textit{\small #2} \\
      \textit{\small #3} & \textit{\small #4} \\
    \end{tabular*}\vspace{-5pt}
}

\begin{document}

% --- EN-TÊTE ---
\begin{center}
    {\Large \textbf{Loïc Aron MBASSI EWOLO}} \\ \vspace{4pt}
    \textbf{\large Stage Ingénieur Développement Web - Imagerie Médicale} \\
    \textit{\small Disponible dès \textbf{Avril 2026} pour 6 mois} \\ \vspace{4pt}
    \small
    \faMapMarker\ Cergy, France (Mobile IDF) \quad
    \faPhone\ (+33) 7 59 52 33 98 \quad
    \faEnvelope\ \href{mailto:loicaron.mbassiewolo@ensea.fr}{loicaron.mbassiewolo@ensea.fr} \\ \vspace{2pt}
    \faLinkedin\ \href{https://www.linkedin.com/in/loic-aron-mbassi-ewolo-3942a9246/}{linkedin} \quad
    \faGithub\ \href{https://github.com/Nameless0l}{github} \quad
    \faGlobe\ \href{https://mbassiloic.tech}{mbassiloic.tech}
\end{center}

\vspace{-6pt}

% --- PROFIL ---
\noindent\small{Élève-ingénieur en double diplôme \textbf{ENSEA/Polytechnique Yaoundé} (5\textsuperscript{ème} année), spécialisé en développement web (\textbf{JavaScript/TypeScript}, React) et traitement d'images. Expérience en télémédecine (SOSDOCTA) et analyse d'images médicales (projet ECG). Je recherche un stage chez \textbf{GE HealthCare} pour contribuer au développement d'applications de visualisation d'images mammographiques.}

% --- FORMATION ---
\section{Formation}
\begin{itemize}[leftmargin=*, label=\tiny$\bullet$, nosep]
    \item \textbf{ENSEA (École Nationale Supérieure de l'Électronique et de ses Applications)} \hfill \textit{Cergy, France | 2025 -- Présent} \\
    \small{Ingénieur 2A, Double Diplôme - Systèmes Embarqués, Traitement du Signal et IA.}
    \vspace{2pt}
    \item \textbf{École Polytechnique de Yaoundé (1\textsuperscript{ère} école d'ingénieur CEMAC)} \hfill \textit{Cameroun | 2021 -- 2025} \\
    \small{Diplôme d'Ingénieur de Conception (5\textsuperscript{ème} année) : Génie Informatique, Génie Logiciel, Architecture Web.}
    \item \textbf{Université de Yaoundé I} \hfill \textit{Cameroun | 2020 -- 2022} \\
    \small{Licence Informatique \& Mathématiques.}
\end{itemize}

% --- COMPÉTENCES TECHNIQUES ---
\section{Compétences Techniques}
\begin{itemize}[leftmargin=*, nosep]
    \item \small{\textbf{Web Frontend :} \textbf{JavaScript/TypeScript}, \textbf{React.js}, HTML5, CSS3, Redux, Responsive Design.}
    \item \small{\textbf{Web Backend :} Node.js, Laravel/PHP, \textbf{REST APIs}, MySQL, MongoDB.}
    \item \small{\textbf{Autres Langages :} Python, \textbf{C/C++}, Java.}
    \item \small{\textbf{Traitement d'Images :} OpenCV, TensorFlow/Keras, Computer Vision.}
    \item \small{\textbf{Outils :} \textbf{Git}, Docker, Figma (notions UI/UX), Méthodologie Agile.}
\end{itemize}

% --- EXPÉRIENCES WEB ---
\section{Expériences en Développement Web}
\begin{itemize}[leftmargin=*, label={-}]

    \resumeSubheading
      {Développeur Full Stack - Plateforme de Télémédecine}{Fév. 2022 -- Mars 2024}
      {SOSDOCTA (Santé / Télémédecine)}{Remote (Belgique/Canada)}
      \begin{itemize}[leftmargin=0.4cm, nosep, label=\tiny$\bullet$]
        \item \small{Développement d'une plateforme complète de télémédecine : réservation, paiement, dashboards médecins/patients.}
        \item \small{Technologies : Laravel, PHP, MySQL, \textbf{JavaScript}, OAuth2/JWT.}
        \item \small{Gestion des données médicales sensibles (conformité RGPD). Maintenance sur 2 ans.}
      \end{itemize}
    \vspace{2pt}

    \resumeSubheading
      {Développeur Frontend - Marketplace}{Fév. 2023 -- Mai 2023}
      {TOGETTECH-LLC}{Yaoundé, Cameroun}
      \begin{itemize}[leftmargin=0.4cm, nosep, label=\tiny$\bullet$]
        \item \small{Développement d'une interface marketplace responsive avec \textbf{React.js} et Redux.}
        \item \small{Création de composants réutilisables, optimisation des performances frontend.}
      \end{itemize}
    \vspace{2pt}

    \resumeSubheading
      {Développeur Backend - API Financière}{Oct. 2024 -- Jan. 2025}
      {CSC Fintech}{Yaoundé, Cameroun}
      \begin{itemize}[leftmargin=0.4cm, nosep, label=\tiny$\bullet$]
        \item \small{Conception d'une architecture \textbf{REST API} sécurisée avec gestion des transactions concurrentes.}
        \item \small{Technologies : Laravel, Docker, Redis, MySQL.}
      \end{itemize}

\end{itemize}

% --- PROJETS IMAGERIE ---
\section{Projets en Traitement d'Images \& Santé}
\begin{itemize}[leftmargin=*, label={-}]

\item \textbf{Détection de Cancer du Poumon (IA Médicale)} \hfill \textit{ENSPY | Sept. 2024 -- Juin 2025}
\vspace{-2pt}
\begin{itemize}[leftmargin=0.4cm, nosep, label=\tiny$\bullet$]
    \item \small{Classification d'\textbf{images médicales} : distinction cas bénins/malins.}
    \item \small{Mise en place de pipelines de \textbf{prétraitement d'images} pour améliorer la précision des modèles.}
\end{itemize}
\vspace{3pt}

\item \textbf{Serious Game pour la Formation Médicale} \hfill \textit{Projet Unity | 2024}
\vspace{-2pt}
\begin{itemize}[leftmargin=0.4cm, nosep, label=\tiny$\bullet$]
    \item \small{Conception d'un simulateur de consultation pour aider les internes en médecine.}
    \item \small{Travail sur l'\textbf{interface utilisateur} et les interactions avec des besoins utilisateurs spécifiques.}
\end{itemize}
\vspace{3pt}

\item \textbf{Projet UROP - Analyse d'Images ECG} \hfill \textit{Mentors Google DeepMind | 2024}
\vspace{-2pt}
\begin{itemize}[leftmargin=0.4cm, nosep, label=\tiny$\bullet$]
    \item \small{Détection d'anomalies cardiaques à partir d'\textbf{images d'électrocardiogrammes}.}
    \item \small{Prétraitement d'images, extraction de features, classification avec TensorFlow/Keras.}
\end{itemize}
\vspace{3pt}

\item \textbf{Outil d'Analyse Vidéo avec YOLO} \hfill \textit{Projet Personnel | 2025}
\vspace{-2pt}
\begin{itemize}[leftmargin=0.4cm, nosep, label=\tiny$\bullet$]
    \item \small{Script Python : détection de personnes, extraction de séquences, optimisation mémoire.}
\end{itemize}

\end{itemize}

% --- RECHERCHE ---
\section{Expérience en Recherche}
\begin{itemize}[leftmargin=*, label={-}]
    \resumeSubheading
      {Assistant de Recherche - IA}{Juin 2025 -- Sept. 2025}
      {MILA (Institut Québécois d'IA) - Remote}{Québec, Canada}
      \begin{itemize}[leftmargin=0.4cm, nosep, label=\tiny$\bullet$]
        \item \small{Travaux sur la généralisation des réseaux de neurones sous Python/PyTorch.}
      \end{itemize}
\end{itemize}

% --- CERTIFICATIONS & LANGUES ---
\section{Hackathons, Certifications \& Langues}
\begin{itemize}[leftmargin=*, nosep]
    \item \textbf{Hackathons :} \textbf{1\textsuperscript{ère} Place} IndabaX Africa (IA Santé : nettoyage données, dashboard Streamlit, API Flask).
    \item \textbf{Certifications :} Supervised ML (DeepLearning.AI), Python for Data Science (IBM).
    \item \textbf{Langues :} Français (Natif), \textbf{Anglais} (Professionnel/Technique).
\end{itemize}

\end{document}
